\section{Security Analysis}
\label{sec:security}

\subsection{What is protected}

Sirus removes the original native body of a selected function and replaces it
with an encrypted bytecode program for a per-build interpreter. An attacker who
reads the protected image sees the host image, the patched entry, the VM engine,
and an encrypted payload. The original arithmetic and control flow are not
present as native instructions. The protected entry is still identifiable through
the host symbol table and the \code{E9} jump; the goal is not to hide that a
function is protected, but to make the protected logic expensive to recover.

\subsection{Objective-by-objective assessment}

\paragraph{O1: readable original logic.}
The body is overwritten with deterministic random bytes after the entry patch.
The patch/scrub test checks that the original bytes change and that the body is
absent from the protected file and from a live memory snapshot. Ghidra and
radare2 reduce the host targets to jump or bad-data stubs. The injected segment
contains engine code, not the original function.

\paragraph{O2: operation obscurity.}
The arithmetic is rewritten through exact MBA identities, constants are split
into XOR shares, and the ISA aliases obscure which semantic operation is
executing. The external MBA study shows that linear forms are reducible and that
the non-linear wrapper raises the cost for the tested general simplifier. This is
measured cost asymmetry, not a proof of irreducibility.

\paragraph{O3: build independence.}
A new seed changes the opcode assignment, operand order, register maps, immediate
endianness, state ids, block order, keys, entry shares, MBA shapes, decoys,
handler order, and code/data layout. The differential payload equal-byte ratio is
0.40\%; the full VM segment is 2.43\% equal at equal offsets. The remaining
similarity is the interpreter skeleton and shared helpers. The metric is a
diversity indicator, not a semantic unlinkability proof.

\paragraph{O4: tamper detection.}
The effective payload key mixes the SipHash of the encrypted payload with the
expected value. A modified payload decrypts to invalid structures and returns
\code{BadPayload} or a deterministic wrong value. The binding does not use a
visible comparison branch. It is not a public-key MAC: an attacker who can
rewrite the expected hash and key shares is outside its scope.

\paragraph{O5: host safety.}
The safe runtime contains no unsafe code, no heap allocation, no host memory
access, and no panic path in normal execution. Kani, Miri, libFuzzer, the
malformed-payload suite, and valgrind cover complementary slices. The claim is
memory safety under the tested models and inputs, not a formal proof of the
entire interpreter.

\paragraph{O6: forced-return resistance.}
The protected function returns data that the host consumes downstream. A forced
return that replaces the protected value with a constant changes the observable
output of the test harness. A caller-side patch that ignores the result is
outside the protected function's scope; no local protector can prevent a caller
from discarding a return value.

\paragraph{O7: string hygiene.}
The generated engine has no formatting, panic strings, or libc imports. The
scoped \code{.sirus} string scan is clean. The host image retains its Rust/std
strings and symbols; Sirus does not strip the host.

\subsection{Attack-cost table}

Table~\ref{tab:attacks} groups the attack classes by the control that addresses
them and the residual that remains after the control.

\begin{table}[htbp]
\centering
\small
\caption{Attack classes, controls, and residuals.}
\label{tab:attacks}
\begin{tabularx}{\textwidth}{@{}p{0.20\textwidth}p{0.34\textwidth}X@{}}
\toprule
Attack & Control & Residual \\
\midrule
Static disassembly & Body removed; flattened encrypted bytecode; encrypted decode table & Entry symbols remain; decompiler work remains \\
MBA simplification & Layered exact identities; non-linear wrapper; constant blinding & Linear forms reduce; SiMBA reduces inverse forms; not a proof \\
Dynamic tracing & Byte-at-a-time decode; per-build states; decoys; masked state & DBI can record semantics; no kernel/DBI detection in safe tier \\
Taint analysis & Arguments influence flags, states, and decode geometry & Taint depth not formally proven \\
Symbolic execution & Large random states; encrypted manifest; MBA; fuel bound & Path explosion is probabilistic; bounded timeouts only \\
Memory dumping & No decrypted body cache; masked registers and stack & A live instruction-local value can be observed \\
Forced return & Result consumed as data downstream & Caller-side patch remains possible \\
Differential analysis & Per-build ISA, keys, layout, MBA, decoys & Interpreter skeleton reused; 2.43\% segment equality \\
VM fuzzing & Checked decode and memory; fuel bound & Semantic garbage may produce a value \\
Timing & Optional fixed-pass oblivious mode & Normal profile leaks; no Spectre proof \\
\bottomrule
\end{tabularx}
\end{table}

\subsection{Differential analysis}

The differential harness protects the same binary with two fixed seeds and
compares the resulting images at equal offsets. The payload is 0.40\% equal; the
full VM segment is 2.43\% equal. The differences include the encrypted manifest
and bodies, opcode and alias tables, register maps, state ids, entry shares,
MBA shapes, decoys, handler order, and code/data placement. The similarities are
the interpreter skeleton, shared runtime helpers, and canonical instruction
sequences.

A stronger differential attacker can align functions and compare semantics
rather than bytes. Sirus does not claim resistance to a semantically aware
diffing system that reconstructs the interpreter; it raises the amount of work
required to do so.

\subsection{Integrity and tamper behavior}

The integrity binding is a 64-bit SipHash-derived value mixed into the effective
key. It detects accidental corruption and naive patching. A modified payload
produces invalid decodes or wrong values; the tamper test verifies that the
result changes without a crash or hang. The binding is not a signature: an
attacker who can modify the generated code, expected hash, and key shares
together can produce a consistent alternative. The safe tier relies on the
combination of encryption, integrity, and structural checks, not on a trusted
hardware root.

\subsection{Memory-dump behavior}

The original body is absent from the live process. The runtime does not cache a
decrypted body; instructions are decoded into locals one at a time. Virtual
registers and the private stack are XOR-masked at rest with per-invocation keys.
A passive snapshot therefore sees masked values. A decoded value can exist
briefly in a local while one instruction executes; an instruction-level tracer
is a stronger attacker than the snapshot diagnostic.

\subsection{Timing and side channels}

The default branchy profile is not constant-time. The explicit mode makes the
decoded block sequence and private-memory address sequence independent of the
inputs under the stated pass bound. It does not make the ALU operations,
cache-bank behavior, branch predictor state, or speculative execution
input-independent. The dudect and Cachegrind results are empirical checks, not a
formal side-channel proof.

\subsection{Residual risk}

The main residuals are:

\begin{itemize}
  \item the interpreter skeleton is shared between builds;
  \item the payload binding is 64-bit and not a public-key MAC;
  \item the normal profile leaks timing;
  \item the constant-time mode has a fixed pass bound and does not model
        microarchitectural effects;
  \item the safe tier has no kernel anti-debug or DBI detection;
  \item the host symbol table and entry patch remain visible;
  \item the supported instruction subset excludes many real functions;
  \item no independent audit or formal obfuscation proof exists.
\end{itemize}

These residuals are recorded so that no deployment decision relies on a stronger
claim than the evidence supports.
